\documentclass[11pt]{article}
\usepackage{icmlko}
\begin{document}

\title{어느 접전을 불러도 되나 --- 신호는 진짜인데, 진짜 동률에서는 그 신호도 침묵한다}
\author{월드모델 실험실}
\date{노트 532 · 2026년 8월 4일}
\maketitle

\begin{abstract}
노트 528 이 \emph{실무의 판은 접전}(후보 다섯 중 $1 \cdot 2$ 등 격차 평균
$16.7$ 점)임을 보였고 노트 531 의 축 차이 수는 접전에서 폭이 $+3.8$ \%p 로
평평했다. \textbf{불러도 되는 접전을 가르는 신호}를 찾았다.

\begin{center}
\begin{tabular}{lrrr}
\hline
접전($<20$ 점) 안에서 & 최저 & 최고 & 폭 \\
\hline
축 차이 수(노트 531) & $53.3\%$ & $57.2\%$ & $+3.8$ \%p \\
\textbf{자루 $32$ 개의 합의도} & $50.6\%$ & $62.3\%$ & $\mathbf{+11.7}$ \textbf{\%p} \\
\hline
\end{tabular}
\end{center}

사전 등록 둘(\emph{접전에서도 단조}, \emph{폭 $10$ \%p 이상}) \textbf{다 통과}.
\textbf{그런데 갈라 보니 반쪽이다} --- 실제 격차 $15\!\sim\!20$ 점에서는 폭이
$+10.8$ \%p 인데 $0\!\sim\!5$ 점에서는 $\mathbf{+1.2}$ \textbf{\%p} 다.
\end{abstract}

\begin{figure}[t]
\centering
\includegraphics[width=\textwidth]{figs/f1.pdf}
\caption{\textbf{가} --- 덮음$\cdot$정확도 곡선. \textbf{나} --- 격차 띠별 폭.
\textbf{다} --- 읽기.}
\end{figure}

\section{먼저 내 수치를 깎는다}

같은 실행에서 \emph{예측 차 크기}(여백)도 재 봤고 접전 폭이 $+22.3$ \%p 로
나왔다. \textbf{그 수는 못 쓴다} --- 여백을 도메인별 \emph{최대값}으로
정규화했고 최상위 칸이 세 도메인 $131$ 쌍뿐이었다.

도메인 \emph{안} 백분위로 다시 내면 이렇다.

\begin{center}
\begin{tabular}{lrr}
\hline
여백 상위 & 덮음 & 정확도 \\
\hline
$100\%$ & $100\%$ & $56.1\%$ \\
$50\%$ & $50\%$ & $59.3\%$ \\
$20\%$ & $20\%$ & $62.4\%$ \\
$10\%$ & $10\%$ & $65.2\%$ \\
$5\%$ & $5\%$ & $67.6\%$ \\
\hline
\end{tabular}
\end{center}

\textbf{덮음 없이 적은 정확도는 쓸 수 없다.} 이것이 제품이 실제로 쓰는 표다
--- ``접전의 $20\%$ 만 답하면 $62.4\%$''.

\section{신호는 진짜다}

의심할 만한 것이 하나 있었다 --- 여백이 큰 접전 쌍은 사실 \emph{덜 접전}인
것 아닌가(격차 $19$ 점과 $2$ 점을 다시 고르는 것 아닌가).

\begin{quote}
\textbf{접전 안에서 여백 $\leftrightarrow$ 실제 격차 순위상관은
$\mathbf{+0.024}$ 다.} 거의 $0$ 이다. \textbf{격차를 다시 고르는 것이 아니라
정말로 ``이건 부를 수 있다''를 아는 것이다.}
\end{quote}

\section{그러나 진짜 동률에서는 침묵한다}

격차 띠를 좁게 잘라 그 안에서만 여백 하위 절반과 상위 절반을 견줬다.

\begin{center}
\begin{tabular}{lrrr}
\hline
실제 격차 & 여백 하위 & 여백 상위 & 폭 \\
\hline
$0\!\sim\!5$ 점 & $51.4\%$ & $52.6\%$ & $\mathbf{+1.2}$ \textbf{\%p} \\
$5\!\sim\!10$ & $51.8\%$ & $57.5\%$ & $+5.8$ \%p \\
$10\!\sim\!15$ & $53.7\%$ & $61.8\%$ & $+8.1$ \%p \\
$15\!\sim\!20$ & $55.1\%$ & $65.9\%$ & $+10.8$ \%p \\
\hline
\end{tabular}
\end{center}

\begin{quote}
\textbf{접전은 균질하지 않다.} ``준접전''($15\!\sim\!20$ 점)은 갈리고
\textbf{진짜 동률}($0\!\sim\!5$ 점)은 안 갈린다. 우리가 얻은 신뢰도는
\emph{접전의 바깥쪽}에서만 작동한다.
\end{quote}

이것은 노트 531 의 반복이다 --- 거기서는 축 차이 수가 ``신뢰도가 필요 없는
곳에서만'' 작동했고, 여기서는 여백이 ``접전의 덜 접전인 쪽에서만'' 작동한다.
\textbf{두 번 같은 모양이 나왔다.}

\section*{조항}

\begin{quote}
\textbf{신뢰도 신호는 덮음과 함께 적는다.} ``상위 $x\%$ 만 답하면 $y\%$''
형태가 아니면 제품에서 쓸 수 없고, 최상위 칸의 표본 수를 같이 적지 않으면
\emph{내가 나를 속인다}(여기서 $131$ 쌍이 $+22.3$ \%p 를 만들었다).
\end{quote}

\section*{읽기 --- 로드맵의 근거가 한 겹 더 단단해졌다}

두 노트가 같은 곳을 가리킨다. \textbf{진짜 동률에 가까운 판은 지금 있는
정보로는 안 갈린다} --- 축을 더 세밀히 보든(노트 531), 모델이 얼마나 확신하는지
보든(여기) 마찬가지다.

\begin{quote}
남은 설명은 하나다 --- \textbf{그 판을 가르는 것이 기획서 안에 없다.}
시점 $\cdot$ 경쟁작 $\cdot$ 운영 $\cdot$ 그때 그 사람들이고, 그것은
\textbf{상태 $\cdot$ 개인 $\cdot$ 개입 층}에 있다. 로드맵이 그 순서인 이유가
이제 \emph{두 가지 다른 측정}으로 지지된다.
\end{quote}

\section*{남은 것}

\begin{itemize}
\item \textbf{자루 합의도와 여백은 같은 것을 재나} --- 둘 다 접전에서 비슷한
폭을 냈다. 겹치면 하나만 쓰면 된다.
\item \textbf{덮음 $20\%$ 에서 $62.4\%$ 가 쓸모 있나} --- 제품 쪽 물음이다.
``스무 건 중 넷만 답하고 나머지는 모른다고 한다''가 실무에서 받아들여지는지.
\item \textbf{$0\!\sim\!5$ 점 쌍이 전체의 얼마인가} --- 그만큼이 원리상
못 푸는 몫이다. 아직 안 셌다.
\end{itemize}

\end{document}
